\documentclass[10pt,a4paper]{amsart}
\usepackage{amsmath,amssymb,amsthm,graphicx}
% Retained minimal preamble for the Stacks Project chapter more-algebra.tex (More on Algebra).
% The chapter file pinned-original-source.tex carries only \input{preamble} on its first line, so no
% preamble text lives inside the chapter itself.  The Stacks Project uses one single shared preamble for
% every chapter; the snapshot stages it identically as sources/stacks-chow/preamble.tex and
% sources/stacks-etale-cohomology/preamble.tex.  The block below is copied mechanically, line for line,
% from sources/stacks-chow/preamble.tex lines 14-18 (the Xy-pic setup) and 154-217 (the native theorem
% declarations together with the mathematical macros).  Nothing was added, removed or redefined; the
% wrapper class and the amsmath/amssymb/amsthm/graphicx base packages are supplied by the assembler.
% Chapter: more-algebra.tex -- More on Algebra.
% For commutative diagrams we use Xy-pic
\usepackage[all]{xy}

% We use 2cell for 2-commutative diagrams.
\xyoption{2cell}

% Theorem environments.
%
\theoremstyle{plain}
\newtheorem{theorem}[subsection]{Theorem}
\newtheorem{proposition}[subsection]{Proposition}
\newtheorem{lemma}[subsection]{Lemma}

\theoremstyle{definition}
\newtheorem{definition}[subsection]{Definition}
\newtheorem{example}[subsection]{Example}
\newtheorem{exercise}[subsection]{Exercise}
\newtheorem{situation}[subsection]{Situation}

\theoremstyle{remark}
\newtheorem{remark}[subsection]{Remark}
\newtheorem{remarks}[subsection]{Remarks}

\numberwithin{equation}{subsection}

% Macros
%
\def\lim{\mathop{\mathrm{lim}}\nolimits}
\def\colim{\mathop{\mathrm{colim}}\nolimits}
\def\Spec{\mathop{\mathrm{Spec}}}
\def\Hom{\mathop{\mathrm{Hom}}\nolimits}
\def\Ext{\mathop{\mathrm{Ext}}\nolimits}
\def\SheafHom{\mathop{\mathcal{H}\!\mathit{om}}\nolimits}
\def\SheafExt{\mathop{\mathcal{E}\!\mathit{xt}}\nolimits}
\def\Sch{\mathit{Sch}}
\def\Mor{\mathop{\mathrm{Mor}}\nolimits}
\def\Ob{\mathop{\mathrm{Ob}}\nolimits}
\def\Sh{\mathop{\mathit{Sh}}\nolimits}
\def\NL{\mathop{N\!L}\nolimits}
\def\CH{\mathop{\mathrm{CH}}\nolimits}
\def\proetale{{pro\text{-}\acute{e}tale}}
\def\etale{{\acute{e}tale}}
\def\QCoh{\mathit{QCoh}}
\def\Ker{\mathop{\mathrm{Ker}}}
\def\Im{\mathop{\mathrm{Im}}}
\def\Coker{\mathop{\mathrm{Coker}}}
\def\Coim{\mathop{\mathrm{Coim}}}

% Boxtimes
%
\DeclareMathSymbol{\boxtimes}{\mathbin}{AMSa}{"02}

%
% Macros for moduli stacks/spaces
%
\def\QCohstack{\mathcal{QC}\!\mathit{oh}}
\def\Cohstack{\mathcal{C}\!\mathit{oh}}
\def\Spacesstack{\mathcal{S}\!\mathit{paces}}
\def\Quotfunctor{\mathrm{Quot}}
\def\Hilbfunctor{\mathrm{Hilb}}
\def\Curvesstack{\mathcal{C}\!\mathit{urves}}
\def\Polarizedstack{\mathcal{P}\!\mathit{olarized}}
\def\Complexesstack{\mathcal{C}\!\mathit{omplexes}}
% \Pic is the operator that assigns to X its picard group, usage \Pic(X)
% \Picardstack_{X/B} denotes the Picard stack of X over B
% \Picardfunctor_{X/B} denotes the Picard functor of X over B
\def\Pic{\mathop{\mathrm{Pic}}\nolimits}
\def\Picardstack{\mathcal{P}\!\mathit{ic}}
\def\Picardfunctor{\mathrm{Pic}}
\def\Deformationcategory{\mathcal{D}\!\mathit{ef}}
\usepackage{cleveref}
\crefname{originallemma}{original lemma}{original lemmas}
\Crefname{originallemma}{Original lemma}{Original lemmas}
\makeatletter
\newlabel{more-algebra-lemma-P-local}{{0BK8}{source}{}{sourcepointer-1}{}}
\newlabel{more-algebra-lemma-P-local@cref}{{[originallemma][0][]0BK8}{[1][0][]source}{}{}{}}
\makeatother
\crefname{originallemma}{original lemma}{original lemmas}
\Crefname{originallemma}{Original lemma}{Original lemmas}
\makeatletter
\newlabel{algebra-remark-fundamental-diagram}{{00E6}{source}{}{sourcepointer-2}{}}
\newlabel{algebra-remark-fundamental-diagram@cref}{{[originallemma][0][]00E6}{[1][0][]source}{}{}{}}
\makeatother
\title{Detecting a property of the fibres on localisations}
\author{The Stacks Project Authors}
\date{stacks-project:more-algebra.tex:local-snapshot-20261012; GNU FDL 1.2 or later}
\begin{document}
\maketitle
\setcounter{subsection}{1}
\noindent\textit{Editorial note.} All delivered mathematical text is the author's own, copied from the pinned Stacks Project chapter file more-algebra.tex; no statement or proof step was rewritten, reordered or removed.\par\medskip
\noindent\textit{Editorial note.} The chapter file pinned-original-source.tex carries only \textbackslash{}input\{preamble\} on its first line, so no preamble text lives inside the chapter. The retained file sources/stacks-more-algebra/retained-preamble.tex is a mechanical, line-for-line copy of the project's single shared preamble as staged in the snapshot (Xy-pic setup plus theorem declarations and mathematical macros). Nothing was added, removed or redefined, and the excerpt serves only as preamble, never as a span coordinate.\par\medskip
\noindent\textit{Editorial note.} Both counted spans start at the \textbackslash{}begin line of their environment, so the statement and the author's proof are delivered whole; surrounding prose of the section is deliberately outside the counted spans and is carried as standard prerequisites instead.\par\medskip
\noindent\textit{Editorial note.} The citation fields (source\_id, source\_version, authors, source\_url, license.name, source\_title) are identical to those of every other card this lane has produced from this chapter, so the batch selector treats them as one source.\par\medskip
\noindent\textit{Allowed standard prerequisites (editorial).}\begin{itemize}
\item Property (B): a property of Noetherian rings that can be checked on local rings. assumed in the statement and used at proof line 12895, source lines 12820-12869. Justifies replacing the verification on fibres by a verification on their local rings.
\item The fibre of a ring map over a prime, identified with the spectrum of the base change to the residue field. used at proof lines 12884-12890 with its description of the primes lying over the given one. Supplies the object whose localisation is computed.
\item Identification of the localisation of a base changed algebra with the localisation at a chosen maximal ideal. used at proof lines 12890-12899. Is the computation that reduces conditions (1) and (3) to each other.
\end{itemize}\medskip
\ifdefined\hypertarget\hypertarget{sourcepointer-1}{}\fi
\noindent\textit{Source reference.} Original lemma 0BK8: focused-seeds/AGP/sources/stacks-snapshot-20261011/tags.txt, source lines 14551--14551. The lemma's own canonical tag, quoted so that readers can resolve the label inside the delivered text; it points at the counted statement itself and carries no proof content.\par\medskip
\ifdefined\hypertarget\hypertarget{sourcepointer-2}{}\fi
\noindent\textit{Source reference.} Original lemma 00E6: focused-seeds/AGP/sources/stacks-snapshot-20261011/tags.txt, source lines 523--523. Navigation pointer to a statement the counted proof refers to but does not contain; nothing load-bearing is delegated to it.\par\medskip
\begin{lemma}
\label{lemma-P-local}
Let $R \to \Lambda$ be a homomorphism of Noetherian rings.
Assume $P$ has property (B). The following are equivalent
\begin{enumerate}
\item the fibres of $R \to \Lambda$ have $P$,
\item the fibres of $R_\mathfrak p \to \Lambda_\mathfrak q$ have $P$
for all $\mathfrak q \subset \Lambda$ lying over $\mathfrak p \subset R$, and
\item the fibres of $R_\mathfrak m \to \Lambda_{\mathfrak m'}$ have $P$
for all maximal ideals $\mathfrak m' \subset \Lambda$
lying over $\mathfrak m$ in $R$.
\end{enumerate}
\end{lemma}

\begin{proof}
Let $\mathfrak p \subset R$ be a prime. Then the fibre over
$\mathfrak p$ is the ring $\Lambda \otimes_R \kappa(\mathfrak p)$
whose spectrum maps bijectively onto the subset of $\Spec(\Lambda)$
consisting of primes $\mathfrak q$ lying over $\mathfrak p$, see
Algebra, Remark \ref{algebra-remark-fundamental-diagram}.
For such a prime $\mathfrak q$ choose a maximal
ideal $\mathfrak q \subset \mathfrak m'$ and set
$\mathfrak m = R \cap \mathfrak m'$.
Then $\mathfrak p \subset \mathfrak m$ and we have
$$
(\Lambda \otimes_R \kappa(\mathfrak p))_\mathfrak q \cong
(\Lambda_{\mathfrak m'} \otimes_{R_\mathfrak m}
\kappa(\mathfrak p))_\mathfrak q
$$
as $\kappa(\mathfrak q)$-algebras. Thus (1), (2), and (3) are equivalent
because by (B) we can check property $P$ on local rings.
\end{proof}

\end{document}
